\section{Introducción}

El fenómeno El Niño – Oscilación del Sur constituye un ejemplo de sistema dinámico acoplado en el que una perturbación térmica del océano puede producir respuestas posteriores en la humedad atmosférica y en la precipitación. En este trabajo práctico se utilizará un modelo temporal simplificado de tres variables para estudiar, por un lado, propiedades numéricas del método de Euler aplicado a un problema de valores iniciales y, por otro, la respuesta del modelo frente a cambios en sus parámetros.

El objetivo es reproducir de manera aproximada un evento climático real, utilizando un sistema de ecuaciones diferenciales ordinarias no lineales con excitación estacional como caso de estudio para analizar convergencia, orden, estabilidad y sensibilidad paramétrica. 
